% =====================================================================
\chapter{Feature Store and Feature Catalogue}
\label{ch:features}
% =====================================================================

\section{Why a feature store and not a script}

The most common cause of a backtest that works and a live system that does not
is \term{train/serve skew}: the offline feature was computed differently from
the online one. The rule is absolute.

\begin{keybox}[title={One definition, two materialisations}]
A feature is defined \emph{once}, in one function, with an explicit dependency
list and an as-of timestamp. The offline path runs it over Parquet history; the
online path runs the same function over Redis state. A CI test replays a live
day offline and asserts every feature matches to within floating-point
tolerance. If it does not match, the build fails.
\end{keybox}

\begin{lstlisting}[language=Python,caption={Feature definition pattern}]
@feature(name="ofi_60s", version=3, deps=["book.bid","book.ask"],
         window="60s", ttl="5s", fill="ffill_max_2")
def order_flow_imbalance(book: BookWindow) -> float:
    """Signed order-flow imbalance over the trailing 60 seconds.
    Point-in-time safe: reads only data with ts <= as_of."""
    ...

# Offline: engine.materialise(ofi_60s, start, end) -> Parquet column
# Online : engine.serve(ofi_60s, as_of=now)        -> float from Redis
# CI     : assert_offline_online_parity(ofi_60s, sample_day)
\end{lstlisting}

\section{Feature catalogue}

\subsection{Group A --- Price and trend (the baseline)}

\begin{longtable}{@{}p{4.6cm}p{4.0cm}p{5.6cm}@{}}
\toprule
\textbf{Feature} & \textbf{Windows} & \textbf{Rationale}\\
\midrule
Log returns & 1, 2, 5, 15, 30, 60 min & Raw momentum at several scales.\\
Return z-score vs time-of-day baseline & same & Normalises the U-shaped
intraday volatility profile. \textbf{Critical.}\\
Realised volatility (Parkinson, Garman--Klass, Rogers--Satchell) & 5, 15, 30, 60
min & Range-based estimators are far more efficient than close-to-close.\\
Bipower variation and jump component & 30 min & Separates continuous
volatility from jumps; events cause jumps.\\
Distance from VWAP, from opening range, from previous close & session & Standard
intraday anchors that real traders actually watch.\\
Position within the day's range & session & Mean-reversion versus breakout
context.\\
Trend strength / efficiency ratio & 15, 30 min & Distinguishes a directional
move from chop; regime information.\\
Autocorrelation of 1-min returns & 30 min & Negative values signal a
mean-reverting microstructure regime.\\
\bottomrule
\end{longtable}

\subsection{Group B --- Microstructure (where the short-horizon edge lives)}

\begin{longtable}{@{}p{4.6cm}p{4.0cm}p{5.6cm}@{}}
\toprule
\textbf{Feature} & \textbf{Windows} & \textbf{Rationale}\\
\midrule
Order-flow imbalance (OFI) & 10, 30, 60 s & Best-in-class short-horizon
directional predictor when book data is available.\\
Bid--ask spread and its z-score & 30 s & Widening spreads precede volatility and
raise transaction cost --- feeds both the model and the cost model.\\
Depth imbalance at top levels & instantaneous & Directional pressure.\\
Trade-sign imbalance (tick rule / Lee--Ready) & 60 s & Aggressor direction.\\
Volume-clock and VPIN-style toxicity proxy & 50 buckets & Flow toxicity spikes
before large moves.\\
Large-trade ratio & 60 s & Institutional participation proxy.\\
Quote update intensity & 10 s & Attention and uncertainty proxy.\\
Futures basis (future minus spot) & instantaneous & Basis moves often lead spot.\\
\bottomrule
\end{longtable}

\begin{warnbox}[title={Availability caveat}]
Full order-book depth may not be available on a retail broker feed. Design the
model so Group~B is optional: train two variants, with and without
microstructure, and let the serving layer select by feed capability. Never
train on features you will not have at 09:16 on a live Tuesday.
\end{warnbox}

\subsection{Group C --- Options-derived (the most differentiated group)}

\begin{longtable}{@{}p{4.8cm}p{3.4cm}p{6.0cm}@{}}
\toprule
\textbf{Feature} & \textbf{Update} & \textbf{Rationale}\\
\midrule
ATM implied volatility, both expiries & 1--5 s & Level of expected volatility.\\
IV term structure slope (near vs next expiry) & 5 s & Inversion signals expected
near-term event risk.\\
25-delta risk reversal (skew) & 5 s & The market's directional fear;
one of the strongest single option-derived predictors.\\
Butterfly / smile convexity & 5 s & Tail pricing.\\
IV rank and IV percentile & daily & Is volatility cheap or expensive in context?\\
Put-call ratio, by volume and by open interest & 1 min & Positioning and
sentiment.\\
Open-interest change by strike, call and put & 1 min & Where new positions are
being built --- writers define support and resistance.\\
Max pain and distance to it & 5 min & Expiry-week pinning effect.\\
Estimated dealer gamma exposure (GEX) & 1 min & Positive gamma dampens moves,
negative gamma amplifies them. \textbf{Regime-defining.}\\
Implied move for the remaining session & 1 min & The denominator of the trade
filter in Chapter~\ref{ch:reality}.\\
Volatility risk premium (IV minus forecast realised) & 5 min & The core edge for
volatility structures.\\
\bottomrule
\end{longtable}

\subsection{Group D --- Cross-asset and global}

Overnight and same-session values of: global equity futures, regional indices,
crude, gold, dollar index, USD/INR, US yields, and a global-risk composite
(first principal component of a basket of risk assets). Include both level
changes and z-scores versus their own trailing distributions.

\subsection{Group E --- Event features}

\begin{longtable}{@{}p{4.4cm}p{10.2cm}@{}}
\toprule
\textbf{Feature} & \textbf{Note}\\
\midrule
Event class one-hot / embedding & Learned embedding is better once enough
history exists.\\
Calibrated severity, polarity, novelty, credibility, confidence & From
Chapter~\ref{ch:event}.\\
Surprise $z$-score & For scheduled releases; the single most powerful macro
feature.\\
Prior impact posterior (mean, spread, $n_{\text{eff}}$) & From
Chapter~\ref{ch:impact}. Feeding the prior in as a feature lets the model learn
when to trust it.\\
Seconds since event & The decay clock.\\
Event density in trailing 30 min & Multiple simultaneous events change dynamics.\\
Conflict flag & Opposing events present.\\
\bottomrule
\end{longtable}

\subsection{Group F --- Calendar, regime and session}

Minutes since open, minutes to close, day of week, days to weekly expiry, days
to monthly expiry, expiry-day flag, budget/policy-day flag, holiday-adjacent
flag, current regime label from the regime model, trailing 20-day realised
volatility percentile, and recent trend regime.

\begin{keybox}[title={Feature hygiene}]
Target roughly 120--250 engineered features, not thousands. Run
\textbf{Purged Group Time-Series CV} feature importance and drop anything that
does not survive. High-cardinality raw features and features with any dependence
on future information are removed automatically by a leakage linter that checks
each feature's dependency timestamps.
\end{keybox}

% =====================================================================
\chapter{Prediction Models}
\label{ch:ml}
% =====================================================================

\section{Labelling: the decision that determines everything}

Predicting ``next candle up or down'' is the wrong label because it does not
correspond to a tradeable decision. Use the \term{triple-barrier} method.

\begin{center}
\begin{tikzpicture}[scale=1.0]
\draw[->,thick,msgrey] (0,0) -- (8.6,0) node[right,font=\scriptsize]{time};
\draw[->,thick,msgrey] (0,-1.9) -- (0,1.9) node[above,font=\scriptsize]{price};
\draw[dashed,msgreen,thick] (0,1.3) -- (7.2,1.3) node[right,font=\scriptsize,black]{upper barrier $+k\sigma$};
\draw[dashed,msred,thick]   (0,-1.3) -- (7.2,-1.3) node[right,font=\scriptsize,black]{lower barrier $-k\sigma$};
\draw[dashed,msblue,thick]  (7.2,-1.7) -- (7.2,1.7) node[above,font=\scriptsize,black]{time barrier $h$};
\draw[msblue,very thick] (0,0) .. controls (1.4,0.4) and (2.4,-0.5) .. (3.2,0.35)
                         .. controls (4.2,1.0) and (4.6,0.9) .. (5.0,1.31);
\fill[msgreen] (5.0,1.31) circle (2.4pt) node[above right,font=\scriptsize,black]{label $= +1$};
\node[font=\scriptsize,anchor=north] at (0,-0.1) {$t_0$ (event)};
\end{tikzpicture}
\end{center}

\begin{longtable}{@{}p{3.6cm}p{11.0cm}@{}}
\toprule
\textbf{Choice} & \textbf{Specification}\\
\midrule
Barriers & Upper and lower at $\pm k\sigma_t$ where $\sigma_t$ is the
\emph{current} volatility estimate (not a fixed point value). $k$ tuned per
horizon; start with $k=1.5$.\\
Time barrier & $h \in \{5, 15, 30\}$ minutes and always before the session close.\\
Label & $+1$ upper hit first, $-1$ lower hit first, $0$ time barrier hit.\\
Sample weights & Weight each sample by uniqueness (overlapping outcome windows
are not independent) and by absolute realised return.\\
Meta-labelling & A second model predicts \emph{whether the primary model's call
should be taken}. Primary model decides side; meta model decides size, including
size zero. This consistently improves precision.\\
\bottomrule
\end{longtable}

\begin{keybox}[title={Why triple-barrier changes the product}]
It directly encodes stop-loss and target into the label. A model trained this way
is answering the question you actually care about --- ``if I enter now with this
stop and this target, do I win?'' --- rather than an academic question about
candle colour.
\end{keybox}

\section{Cross-validation without lying to yourself}

\begin{longtable}{@{}p{3.8cm}p{10.8cm}@{}}
\toprule
\textbf{Technique} & \textbf{Why it is mandatory}\\
\midrule
Purged K-fold & Remove training samples whose outcome window overlaps the test
window. Without purging, the model sees the answer.\\
Embargo & Additionally drop a buffer (e.g. 1 trading day) after each test fold
to kill serial-correlation leakage.\\
Walk-forward (anchored and rolling) & The only evaluation that resembles
production. Report both.\\
Combinatorial purged CV & Generates many backtest paths, giving a
\emph{distribution} of Sharpe ratios rather than one lucky number.\\
Deflated Sharpe ratio & Corrects for the number of strategy variants tried.
If you tested 200 configurations, your best Sharpe is inflated --- quantify it.\\
Event-level splitting & Never split the same event's samples across train and
test. Group by event ID.\\
\bottomrule
\end{longtable}

\begin{warnbox}[title={Leakage checklist --- run before believing any result}]
\begin{enumerate}
\item Does any feature use data with a timestamp after the decision time?
\item Is normalisation (mean/std, scaler) fitted on the full dataset instead of
train only?
\item Are volatility estimates computed with a centred window?
\item Are labels from overlapping windows present in both train and test?
\item Was feature selection done before the split?
\item Does the news timestamp reflect publication or the time you downloaded it?
\item Is the option chain snapshot the one available at decision time, or the
end-of-day file?
\item Are corporate actions and index reconstitutions handled point-in-time?
\end{enumerate}
Any single ``yes'' invalidates the result. Most spectacular backtests die here.
\end{warnbox}

\section{Model portfolio}

Do not build one model. Build a small ensemble of models with genuinely
different inductive biases, then combine them.

\begin{longtable}{@{}p{3.2cm}p{3.4cm}p{7.8cm}@{}}
\toprule
\textbf{Model} & \textbf{Predicts} & \textbf{Role and notes}\\
\midrule
M0 Rules baseline & direction & Simple event-prior lookup plus momentum filter.
\textbf{Every ML model must beat this or it does not ship.} Also the
degraded-mode fallback.\\
M1 LightGBM & $P(\text{up})$, triple-barrier class & The workhorse. Handles
tabular heterogeneous features and missing values, trains in minutes, is
interpretable via SHAP. Expect this to be your best model for a long time.\\
M2 Quantile LightGBM & $q_{05},q_{25},q_{50},q_{75},q_{95}$ of $r_h$ &
Pinball loss. Feeds the fan chart directly.\\
M3 Temporal CNN / TCN & direction and magnitude & Learns local shape patterns
from raw candle and microstructure sequences that hand-crafted features miss.\\
M4 Sequence transformer & multi-horizon distribution & Attention over the last
$N$ minutes plus event tokens. Highest ceiling, highest overfitting risk, needs
the most data. \textbf{Build this last, not first.}\\
M5 Volatility model & $\sigma_{t+h}$ & HAR-RV baseline plus a gradient-boosting
correction on the residual. Volatility is the most forecastable target ---
give it a dedicated, well-tuned model.\\
M6 IV-change model & $\Delta\sigma_{IV}$ & Predicts implied volatility movement.
\textbf{Directly determines options profitability.} Features: current IV rank,
event class, term-structure slope, time to expiry, realised-implied spread.\\
M7 Meta-labeller & take / skip & Trained on primary model outputs plus market
state. Its job is to say no.\\
M8 Regime classifier & regime label & Hidden Markov model or clustering over
volatility, trend and liquidity. Everything else is conditioned on it.\\
\bottomrule
\end{longtable}

\section{Ensembling and calibration}

\begin{enumerate}
\item \term{Stacking with time-aware folds.} Out-of-fold predictions from
M1--M4 become inputs to a simple logistic regression meta-learner. Keep the
meta-learner simple; a complex stacker overfits immediately.
\item \term{Regime-conditional weights.} Model weights vary by regime from M8;
a transformer that shines in trending regimes may be useless in chop.
\item \term{Probability calibration.} Fit isotonic regression or Platt scaling on
a held-out, chronologically later slice. Then verify with a reliability diagram
and the Brier score decomposition (reliability, resolution, uncertainty).
\item \term{Conformal prediction.} Wrap the magnitude forecast in split conformal
intervals to obtain finite-sample coverage guarantees. When the interval is too
wide to be tradeable, the honest answer is no trade --- and now the system can
prove it.
\end{enumerate}

\begin{keybox}[title={The calibration contract}]
The dashboard's headline metric is not accuracy. It is the reliability diagram:
of all the times the system said 60\%, how often was it right? A system that is
calibrated can be sized correctly. A system that is accurate but overconfident
will still blow up.
\end{keybox}

\section{Training loop and retraining policy}

\begin{longtable}{@{}p{3.6cm}p{3.2cm}p{7.6cm}@{}}
\toprule
\textbf{Job} & \textbf{Cadence} & \textbf{Detail}\\
\midrule
Nightly feature materialisation & Daily 18:00 & Recompute the day's features,
run quality gates, append to Parquet.\\
Incremental retrain & Weekly (weekend) & Refit M1, M2, M5, M6 on the expanded
window; run walk-forward; promote only if it beats the incumbent on the
last 4 weeks out-of-sample.\\
Full retrain and architecture search & Quarterly & Includes M3/M4
hyperparameters and feature re-selection.\\
Impact library rebuild & Weekly & New events added; posteriors updated.\\
Calibration refresh & Weekly & Isotonic refit on the most recent 4 weeks.\\
Drift monitor & Continuous & PSI on features, KL divergence on prediction
distribution, rolling Brier score. Breach triggers alert and blocks
auto-execution.\\
Champion/challenger & Continuous & The challenger runs in shadow mode on live
data for at least 15 sessions before it may be promoted.\\
\bottomrule
\end{longtable}

\begin{warnbox}[title={Anti-overfitting discipline}]
Keep a \textbf{locked holdout}: the most recent 3 months of data, encrypted and
untouched. You may evaluate against it \emph{once}, at the Phase~7 gate.
Every additional peek burns the only honest estimate you will ever have.
Log every experiment; deflate your Sharpe by the number of trials.
\end{warnbox}
